\begin{abstract}
Understanding AI tutors requires knowing both whether they answer correctly
and how they usually help students. We study \emph{educational personality}:
recurring, observable teaching tendencies. Our investigation begins with
selected subsets of over one million educational evaluation records and
continues with targeted experiments and tests on new responses. Models show
recurring differences in how directly they help and how much information
they present, with stronger evidence for directness. Other measures show
shared responses or differences limited to particular settings. Knowing how
often a model previously gave answers or invited student reasoning helps
predict these actions on new problems, including under tested changes in
student wording. For acknowledgement of feeling or difficulty, knowing how
each model responds to student cues improves prediction beyond its usual
rate, but this extra benefit is lost under new wording. Explicit teaching
instructions make models perform the same required action while leaving
other choices different. These findings help educators understand and compare
AI tutors through their teaching habits, when those habits remain predictable,
and how they change with instructions. The findings also help developers investigate model selection, prompt design,
and additional system support.
\end{abstract}
